\documentclass[12 pt, letterpaper]{hmcpset}
\usepackage{hw-preamble}

\name{Jayden Thadani}
\class{MATH 142 --- Differential Geometry}
\assignment{Homework assignment \#8}
\duedate{29th March 2024}

\begin{document}

\begin{problem}[Problem 1.]
	\begin{enumerate}
		 \item Say $F, G : \mathbb{R}^{3} \to \mathbb{R}^{3}$ are isometries. Show that $\operatorname{sgn}(F \circ G) = \operatorname{sgn}(F) \operatorname{sgn}(G)$.
		\item Show that $F^{-1} : \mathbb{R}^{3} \to \mathbb{R}^{3}$ and $F$ have the same sign. That is, $\operatorname{sgn}(F) = \operatorname{sgn}(F^{-1})$.
	\end{enumerate}
\end{problem}

\begin{solution}
	\begin{enumerate}
		\item Suppose $F = T \circ C$ and $G = S \circ D$ where $T, S$ are translations $T : \mathbf{p} \mapsto \mathbf{p} + \mathbf{a}$ and $S : \mathbf{p} \mapsto \mathbf{p}$ and $C, D$ are orthogonal transformations. 

			Then since $C$ is a linear map (all orthogonal transformations are linear), we have
			\[
				\begin{split}
					(F \circ G) (\mathbf{p}) & = F((S \circ D)(\mathbf{p})) \\
								 & = F(D(\mathbf{p}) + \mathbf{b}) \\
								 & = (T \circ C)(D(\mathbf{p}) + \mathbf{b}) \\
								 & = C(D(\mathbf{p}) + \mathbf{b}) + \mathbf{a} \\
								 & = (C \circ D)(\mathbf{p}) + (C(\mathbf{b}) + \mathbf{a}).
				\end{split}
			\]
			Thus, we can write $F \circ G$ as the composition $R \circ E$ where $E = C \circ D$ is an orthogonal transformation, and $R$ is translation by $C(\mathbf{b}) + \mathbf{a}$. Thus,
			\[
				\begin{split}
					\operatorname{sgn}(F \circ G) & = \det E \\
								      & = \det (C \circ D) \\
								      & = (\det C) (\det D) \\
								      & = \operatorname{sgn} (F) \operatorname{sgn} (G)
				\end{split}
			\]
			as desired, where the third equality is due to the multiplicity of the determinant for linear operators.
		\item We know that $F \circ F^{-1} = I$ where $I$ is the identity map on $\mathbb{R}^{3}$. Since $I$ is an orthogonal map itself, we have
			\[
				\begin{split}
					\operatorname{sgn} I & = \det I \\
							     & = 1.
				\end{split}
			\]
			Thus,
			\[
				\operatorname{sgn}(F) \operatorname{sgn}(F^{-1}) = 1.
			\] 

			Since their product is positive, they must have the sign. Since we already know that the sign of any isometry has magnitude $1$, they must be equal.
	\end{enumerate}
\end{solution}

\pagebreak

\begin{problem}[Problem 2.]
	Recall that any orthogonal transformation $C : \mathbb{R}^{3} \to \mathbb{R}^{3}$ can be uniquely expressed in the form $C (\mathbf{x}) = A \mathbf{x}$ for some $3 \times 3$ orthogonal matrix $A$; we define $\operatorname{sgn}(C) = \det A$. A \textit{rotation} is an is an orthogonal transformation $C$  such that $\operatorname{sgn}(C) = 1$ .
	\begin{enumerate}
		\item Show that $A$ has eigenvalue $1$ as well as a unit eigenvector $\mathbf{e}_{3}$ satisfying $A \mathbf{e}_{3} = \mathbf{e}_{3}$.
		\item Prove that $C$ does, in fact, rotate $\mathbb{R}^{3}$ about an axis $L$. Explicitly, given a rotation $C$, show that there exists a number $\vartheta$ and points $\mathbf{e}_{1}, \mathbf{e}_{2}, \mathbf{e}_{3} \in \mathbb{R}^{3}$ with
			\[
				\mathbf{e}_{i} \cdot \mathbf{e}_{j} = \begin{cases}
					1 & \text{if } i = j \text{, and} \\
					0 & \text{otherwise}
				\end{cases}
			\]
			such that
			\[
				\begin{split}
					C(\mathbf{e}_{1}) & = \cos{\vartheta} \, \mathbf{e}_{1} + \sin{\vartheta} \, \mathbf{e}_{2} \\
					C(\mathbf{e}_{2}) & = - \sin{\vartheta} \, \mathbf{e}_{1} + \cos{\vartheta} \, \mathbf{e}_{2} \\
					C(\mathbf{e}_{3}) & = \mathbf{e}_{3}.
				\end{split}
			\]
	\end{enumerate}
\end{problem}

\begin{solution}
	\begin{enumerate}
		\item Since $C$ is an orthogonal transformation, it preserves norms. Thus, if $\mathbf{v}$ is an eigenvector of $C$ with eigenvalue $\lambda$, we must have
			\[
				\begin{split}
					& \lVert C(\mathbf{v}) \rVert = \lVert \mathbf{v} \rVert \\
					\implies & \lVert \lambda \rVert \lVert \mathbf{v} \rVert = \lVert \mathbf{v} \rVert \\
					\implies & \lvert \lambda \rvert \lVert \mathbf{v} \rVert = \lVert \mathbf{v} \rVert \\
					\implies & \lvert \lambda \rvert = 1. 
				\end{split}
			\]
			Since $C$ must have three eigenvalues counting multiplicity, and for real matrices $\lambda$ is an eigenvalue if and only if its complex conjugate $\overline{\lambda}$ is, we can write the (possibly complex) eigenvalues of $C$ as $\lambda_{1}, \lambda_{2}, \lambda_{3}$. 

			If all of them are real, we must have at least one of them be positive since
			\[
				\det A = \lambda_{1} \lambda_{2} \lambda_{3}.
			\]
			If we had $\lambda_{1} = \lambda_{2} = \lambda_{3} = -1$, then we would have $\det A = -1$, which we know it is not. Thus we can assume without loss of generality that $\lambda_{3} = 1$. 

			If one of them is complex (say without loss of generality $\lambda_{1} \in \mathbb{C} \setminus \mathbb{R}$) then since complex eigenvalues of a real matrix come in conjugate pairs, we must have $\overline{\lambda_{1}}$ as an eigenvalue as well. Suppose without loss of generality that $\lambda_{2} = \overline{\lambda_{1}}$. Then
			\[
				\begin{split}
					\lambda_{3} & = \frac{\det A}{\lambda_{1} \lambda_{2}} \\
						    & = \frac{1}{\lambda_{1} \overline{\lambda_{1}}} \\
						    & = \frac{1}{\lvert \lambda_{1} \rvert^{2}} \\
						    & = 1.
				\end{split}
			\]
			Thus, we have $\lambda_{3} = 1$, an eigenvalue of $C$. 

			Since $\lambda_{3}$ is an eigenvalue, there exists some nonzero $\mathbf{v} \in \mathbb{R}^{3}$ such that $C(\mathbf{v}) = \mathbf{v}$. If we choose $\mathbf{e}_{3} = \mathbf{v}/\lVert \mathbf{v} \rVert$, we get a unit eigenvector with eigenvalue $\lambda_{3} = 1$ as desired.
		\item We choose $\mathbf{e}_{3}$ from the previous part to be a unit eigenvector of eigenvalue $1$. Note that $C(\mathbf{e}_{3}) = \mathbf{e}_{3}$ as desired. Also note that since $\mathbf{e}_{3}$ is a unit vector we have $\mathbf{e}_{3} \cdot \mathbf{e}_{3} = 1$, satisfying the orthonormal basis condition for $i = j = 3$. 

			Let $U = \operatorname{span} \mathbf{e}_{3}$ be the line spanned by $\mathbf{e}_{3}$. Then we can consider its orthogonal complement $U^{\perp}$. Since $\operatorname{dim} U^{\perp} = \operatorname{dim} \mathbb{R}^{3} - \operatorname{dim} U$, it has an orthonormal basis of two vectors $\mathbf{e}_{1}$ and $\mathbf{e}_{2}$. Since they form an orthonormal basis, the condition is satisfied for $i, j \in \{ 1, 2 \}$. Finally, since they are in $U^{\perp}$ they are orthogonal to $\mathbf{e}_{3} \in U$, satisfying the remaining orthonormality conditions. That is, $\mathbf{e}_{1}, \mathbf{e}_{2}, \mathbf{e}_{3}$ is an orthonormal basis. 

			If we write the matrix of $C$ with respect to the basis $\mathbf{e}_{1}, \mathbf{e}_{2}, \mathbf{e}_{3}$, we get
			\[
				\begin{bmatrix}
					a & b & 0 \\
					c & d & 0 \\
					0 & 0 & 1
				\end{bmatrix}
			\]
			where $C(\mathbf{e}_{1}), C(\mathbf{e}_{2})$ remain in $\operatorname{span}(\mathbf{e}, \mathbf{e}_{2})$ since they must remain orthonormal to $\mathbf{e}_{3}$ (orthonormal transformations preserve dot products). Additionally, to preserve orthonormality, we must have the column vectors of the matrix form an orthonormal list, and the determinant of the matrix must be $1$. This gives us the following system of equations
			\[
				\begin{aligned}
					a^{2} + c^{2} & = 1 & b^{2} + d^{2} & = 1 \\
					ab + cd & = 0 & ad - bc & = 1.
				\end{aligned}
			\]
			Since $a^{2} + c^{2} = b^{2} + d^{2} = 1$, we can always find $\vartheta$ such that $a^{2} = d^{2} = \cos^{2}{\vartheta}$ and $b^{2} = c^{2} = \sin^{2}{\vartheta}$. In order to satisfy both of the equations on the second line, we should choose $a = d = \cos{\vartheta}$ and $-b = c = \sin{\vartheta}$. This choice gives us
			\[
				\begin{split}
					ab + cd & = \cos{\vartheta} (- \sin{\vartheta}) + \sin{\vartheta} \cos{\vartheta} \\
						& = 0.
				\end{split}
			\]
			and
			\[
				\begin{split}
					ad - bc & = \cos{\vartheta} \cos{\vartheta} - (- \sin{\vartheta}) \sin{\vartheta} \\
						& = \cos^{2}{\vartheta} + \sin^{2}{\vartheta} \\
						& = 1.
				\end{split}
			\]
			as desired.
	\end{enumerate}
\end{solution}

\pagebreak

\begin{problem}[Problem 3.]
	Let $\mathbf{a}$  be a point of $\mathbb{R}^{3}$  such that $\lVert \mathbf{a} \rVert = 1$ .
	\begin{enumerate}
		\item Prove that the formula
			\[
				C(\mathbf{p}) = \mathbf{a} \times \mathbf{p} + (\mathbf{p} \cdot \mathbf{a}) \mathbf{a}
			\]
			defines an orthogonal transformation $C : \mathbb{R}^{3} \to \mathbb{R}^{3}$.
		\item Describe its general effect on $\mathbb{R}^{3}$.
	\end{enumerate}
\end{problem}

\begin{solution}
	\begin{enumerate}
		\item First we will show that $C$ is linear.
			\[
				\begin{split}
					C(\alpha \mathbf{p} + \beta \mathbf{q}) & = \mathbf{a} \times (\alpha \mathbf{p} + \beta \mathbf{q}) + ((\alpha \mathbf{p} + \beta \mathbf{q}) \cdot \mathbf{a}) \mathbf{a} \\
								   & = \mathbf{a} \times (\alpha \mathbf{p}) + \mathbf{a} \times (\beta \mathbf{q}) + ((\alpha \mathbf{p}) \cdot \mathbf{a} +  (\beta \mathbf{q}) \cdot \mathbf{a}) \mathbf{a} \\
								   & = \alpha(\mathbf{a} \times \mathbf{p}) + \beta (\mathbf{a} \times \mathbf{q}) + \alpha(\mathbf{p} \cdot \mathbf{a}) \mathbf{a} + \beta (\mathbf{p} \cdot \mathbf{a}) \mathbf{a} \\
								   & = \alpha(\mathbf{a} \times \mathbf{p} + (\mathbf{p} \cdot \mathbf{a}) \mathbf{a}) + \beta (\mathbf{a} \times \mathbf{q} + (\mathbf{q} \cdot \mathbf{a}) \mathbf{a}) \\
								   & = \alpha C(\mathbf{p}) + \beta C(\mathbf{q}).
				\end{split}
			\] 

			We will now show that $C$ preserves norms.
			\[
				\begin{split}
					\lVert C(\mathbf{p}) \rVert^{2} & = \lVert \mathbf{a} \times \mathbf{p} + (\mathbf{p} \cdot \mathbf{a}) \mathbf{a} \rVert^{2} \\
									& = \lVert \mathbf{a} \times \mathbf{p} \rVert^{2} + \lVert (\mathbf{p} \cdot \mathbf{a}) \mathbf{a} \rVert^{2} \\
									& = \lVert \mathbf{a} \rVert^{2} \lVert \mathbf{p} \rVert^{2} - (\mathbf{p} \cdot \mathbf{a})^{2} + (\mathbf{p} \cdot \mathbf{a})^{2} \lVert \mathbf{a} \rVert^{2} \\
									& = \lVert \mathbf{p} \rVert^{2}.
				\end{split}
			\]
			Here, the third step follows by Lagrange's identity, and the last step follows by the given fact that $\lVert \mathbf{a} \rVert = 1$. Since both are positive, we can take the square root on both sides to get $\lVert C(\mathbf{p}) \rVert = \lVert \mathbf{p} \rVert$. 

			Using both of these, we can see that $C$ preserves distances. Specifically, 
			\[
				\begin{split}
					d(C(\mathbf{p}), C(\mathbf{q})) & = \lVert C(\mathbf{p}) - C(\mathbf{q}) \rVert \\
									& = \lVert C(\mathbf{p} - \mathbf{q}) \rVert \\
									& = \lVert \mathbf{p} - \mathbf{q} \rVert \\
									& = d(\mathbf{p}, \mathbf{q}).
				\end{split}
			\] 

			Since $C$ preserves norms and distances, it is an orthogonal transformation.
		\item $C$ rotates any point by an angle of $\pi/2$ about the axis defined by $\operatorname{span}(\mathbf{a})$.
	\end{enumerate}
\end{solution}

\pagebreak

\begin{problem}[Problem 4.]
	Let $F : \mathbb{R}^{3} \to \mathbb{R}^{3}$ be an isometry of $\mathbb{R}^{3}$, and write $F = T \circ C$ for some translation $T$ and some orthogonal transformation $C$. Let $\beta : I \to \mathbb{R}^{3}$ be a unit speed curve in $\mathbb{R}^{3}$. Recall that the Frenet-Serret apparatus for $\overline{\beta} = F \circ \beta$ is
	\begin{align*}
		\overline{\mathbf{T}} & = F_{*} \circ \mathbf{T} & \overline{\mathbf{N}} & = F_{*} \circ \mathbf{N} & \overline{\mathbf{B}} & = \operatorname{sgn} F (F_{*} \circ \mathbf{B}) \\
		\overline{v} & = 1 & \overline{\kappa} & = \kappa & \overline{\tau} & = \operatorname{sgn} F \, \tau.
	\end{align*}
	\begin{enumerate}
		\item Show that if $\beta$ is a cylindrical helix, then $\overline{\beta} = F \circ \beta$ is also a cylindrical helix.
		\item The \textit{spherical image} $\sigma : I \to \mathbb{R}^{3}$ of $\beta$ is that function with the same Euclidean coordinates as that of $\beta'(t)$. That is, $\mathbf{T}(t) = [\sigma(t), \beta(t)]$. Show that, if $\beta$ has spherical image $\sigma$, then $\overline{\beta}$ has spherical image $C \circ \sigma$.
	\end{enumerate}
\end{problem}

\begin{solution}
	\begin{enumerate}
		\item If $\beta$ is a cylindrical helix, then we have some $[\mathbf{u}, \mathbf{p}] \in \mathrm{T}_{\mathbf{p}}\mathbb{R}^{3}$ with fixed $\mathbf{u}$ such that $\mathbf{T}(t) \cdot [\mathbf{u}, \mathbf{p}]$ is constant at each $\mathbf{p} = \beta(t)$. 

			Consider $\overline{[\mathbf{u}, \mathbf{p}]} = F_{*}([\mathbf{u}, \mathbf{p}])$ for the unit tangent vector for $\overline{\beta}$. Since we know that $F_{*}$ preserves dot products, we have
			\[
				\begin{split}
					\overline{\mathbf{T}}(t) \cdot \overline{[\mathbf{u}, \mathbf{p}]} & = F_{*}(\mathbf{T}(t)) \cdot F_{*}(\overline{[\mathbf{u}, \mathbf{p}]}) \\
													   & = \mathbf{T}(t) \cdot [\mathbf{u}, \mathbf{p}]
				\end{split}
			\]
			and thus, $\overline{\mathbf{T}}(t) \cdot \overline{[\mathbf{u}, \mathbf{p}]}$ is constant too. Thus, $\overline{\beta}$ is a cylindrical helix.
		\item We know that
			\[
				\begin{split}
					\overline{\mathbf{T}}(t) & = F_{*}(\mathbf{T}(t)) \\
								 & = F_{*}([\sigma(t), \beta(t)]) \\
								 & = [(C \circ \sigma)(t), F(\beta(t))] \\
								 & = [(C \circ \sigma)(t), \overline{\beta}(t)].
				\end{split}
			\]
			Thus, $\overline{\beta}$ has the spherical image $C \circ \sigma$.
	\end{enumerate}
\end{solution}

\pagebreak

\begin{problem}[Problem 5.]
	Let $F$ be an isometry of $\mathbb{R}^{3}$. For each vector field $\mathbf{V}$ which sends a point $\mathbf{p}$ to the tangent vector $\mathbf{V}(\mathbf{p}) = [\mathbf{v}, \mathbf{p}]$, let $\overline{\mathbf{V}}$ be that vector field which sends a point $\mathbf{q} = F(\mathbf{p})$ to the tangent vector $F_{*}(\mathbf{V}(\mathbf{p}))$. Prove that \textit{isometries preserve covariant derivatives}. That is, show $\overline{\nabla_{\mathbf{V}} \mathbf{W}} = \nabla_{\overline{\mathbf{V}}} \overline{\mathbf{W}}$.
\end{problem}

\begin{solution}
	Let $F = T \circ C$ where $C$ is an orthogonal transformation and $T$ is a translation. Let $\mathbf{W}(\mathbf{p}) = [\mathbf{w}, \mathbf{p}]$ where $\mathbf{w}$ is a function $\mathbf{w} : \mathbb{R}^{3} \to \mathbb{R}^{3}$.

	We know, by definition, that $\nabla_{\mathbf{V}} \mathbf{W}(\mathbf{p}) = [D\mathbf{w} \Big |_{\mathbf{p}}(\mathbf{v}), \mathbf{p}]$. Also, notice that for $\mathbf{q} = F(\mathbf{p})$, we have $\overline{\mathbf{V}}(\mathbf{q}) = [(C \circ \mathbf{v} \circ F^{-1})(\mathbf{q}), \mathbf{q}]$, and similarly for $\overline{\mathbf{W}}$. Thus, for some $\mathbf{q} = F(\mathbf{p})$,
	\[
		\begin{split}
			\nabla_{\overline{\mathbf{V}}} \overline{\mathbf{W}}(\mathbf{q}) & = [D(C \circ \mathbf{w} \circ F^{-1}) \Big |_{\mathbf{p}} (C \circ \mathbf{v})(\mathbf{p}), \mathbf{q}] \\
									      & = \left [ DC \Big |_{\mathbf{w}(\mathbf{p})} \left ( D \mathbf{w} \Big |_{\mathbf{p}} \left ( DF^{-1} \Big |_{\mathbf{q}} (C(\mathbf{v}(\mathbf{p}))) \right ) \right ), \mathbf{q} \right ] \\
									      & = \left [ C \left ( D\mathbf{w} \Big |_{\mathbf{p}} (C^{-1}(C(\mathbf{v}(\mathbf{p})))) \right ), \mathbf{q} \right ] \\
									      & = \left [ C \left ( D \mathbf{w} \Big |_{\mathbf{p}} (\mathbf{v}(\mathbf{p})) \right ), \mathbf{q} \right ] \\
									      & = F_{*} \left ( \left [ D\mathbf{w} \Big |_{\mathbf{p}} (\mathbf{v}(\mathbf{p})), \mathbf{p} \right ] \right ) \\
									      & = F_{*}(\nabla_{\mathbf{V}} \mathbf{W}(\mathbf{p})) \\
									      & = \overline{\nabla_{\mathbf{V}} \mathbf{W}}{(\mathbf{q})}.
		\end{split}
	\]
	Here the second step follows by the multivariable chain rule, the third step follows by noticing that the orthogonal part of $F^{-1}$ is $C^{-1}$ and that the derivative of any linear map is itself, and the steps after those just follow by the definitions. 

	Note that we used the notation $Df \Big |_{\mathbf{a}}$ to denote the linear map that is the total derivative of the function $f$, at the point $\mathbf{a}$, and the notation $Df \Big |_{\mathbf{a}}(\mathbf{v})$ to denote evaluating that linear map on the vector $\mathbf{v}$.
\end{solution}

\end{document}
